%! TeX root = ../phd-thesis.tex
%! Author = davidedomini

%----------------------------------------------------------------------------------------
\chapter{Collective Adaptive Systems}
\label{chap:context}
\minitoc
%----------------------------------------------------------------------------------------

\section{Foundations}
\label{sec:cas-foundations}

\subsection{Definition and Scope}
\label{subsec:cas-definition}

A \glsentryfull{CAS} is a system whose behavior results
 from the interactions among multiple autonomous or semi-autonomous entities 
 and that changes during operation in response to variations in either the system itself 
 or its environment~\cite{DBLP:conf/huc/Ferscha15}. 
% 
The term identifies two complementary properties. 
%
A system is \emph{collective} when its function cannot be attributed to any component in isolation, 
 but depends on how multiple components interact and coordinate. 
% 
It is \emph{adaptive} when it can modify the mechanisms that govern its behavior, 
 rather than merely produce different outputs according to a fixed set of rules. 
% 
Adaptation may affect parameters, decision policies, interaction patterns, role assignments, 
 or the organization of the collective itself.

This definition describes a broad class of systems. 
%
Their components may be software agents, sensors, mobile devices, robots, vehicles, users, or services;
 they may be homogeneous or differ in capabilities and responsibilities; 
 and their interactions may take place in a physical environment, a communication network, or both. 
%
What makes them a collective is not their mere coexistence,
 but the presence of a system-level function or behavior sustained by their interdependence. 
% 
A collection of unrelated devices is therefore not a CAS,
 whereas a population whose members exchange information, divide work, 
 or jointly affect a shared environment can be considered one even when 
 no member has a complete representation of the system.

\glspl{CAS} can be viewed as a manifestation of the transition from computing 
 with individual machines to computing through ensembles of networked entities. 
%
In the collective-computing perspective, 
 computation is no longer confined to a single device or delivered independently by many devices; 
 it is produced by their many-to-many interactions~\cite{DBLP:journals/computer/Abowd}. 
% 
This shift changes the primary object of engineering. 
%
The behavior of each component remains relevant, but it must be understood in relation
 to the organization and dynamics of the ensemble to which it belongs.

\subsection{The Micro--Macro Relationship}
\label{subsec:cas-micro-macro}

The distinction between the \emph{micro} and \emph{macro} levels 
 is central to the characterization of a CAS. 
% 
The micro level comprises the state, observations, decisions, 
 and direct interactions of the individual entities.
%
The macro level concerns properties attributed to the collective as a whole,
 such as spatial coverage, consensus, a partition into functional regions, 
 the motion of a swarm, or the quality of a model learned from distributed data.
%
These properties are not stored or controlled by a single component in the general case. 
%
Instead, 
 they are induced by many local computations 
 and by the information that propagates through the interaction network.

The relationship between the two levels is bidirectional. 
%
Local decisions produce and continuously reshape macroscopic structures, 
 while those structures alter the context in which subsequent local decisions are made.
%
For instance,
 the selection of a coordinator by nearby devices creates a region at the macro level; 
 membership in that region then determines the information available to each device 
 and the agents with which it cooperates. 
% 
Similarly, 
 local placement or routing decisions can create congestion, 
 which in turn changes the cost of future decisions. 
% 
A \gls{CAS} is therefore better understood as a coupled feedback process 
 than as a static collection of components.

System boundaries are also dependent on the level 
 and purpose of the analysis.
%
A collective may contain transient sub-collectives, 
 and the same entity may participate in several groups associated with different functions. 
% 
Conversely,
 an entire collective may act as one component of a larger system. 
% 
This nesting is particularly common in edge--cloud environments, 
 where devices may form local coordination groups while simultaneously interacting with shared edge 
 and cloud infrastructure. 
% 
Consequently, membership, neighborhood, and role should generally be treated 
 as dynamic properties rather than fixed design-time facts.

\subsection{Adaptation, Self-Organization, and Emergence}
\label{subsec:cas-adaptation}

Adaptation in a \gls{CAS} may occur at both the individual and collective levels. 
%
An individual entity may update a model, revise a decision policy, 
 change its operating parameters, or migrate a computation. 
% 
At the collective level, 
 the system may modify its communication topology, redistribute responsibilities,
 form or dissolve groups, elect coordinators, or relocate services. 
% 
These forms of adaptation are often interdependent: 
 a change in an individual model can alter which entities should cooperate, 
 while a change in collective organization determines the observations 
 and resources available to each learning process.

The ability to adapt is commonly expressed through \emph{self-*} properties.
%
Self-configuration concerns the autonomous selection of a suitable configuration; 
 self-optimization aims to improve performance during operation;
 self-healing enables recovery from failures or degraded conditions; and
 self-protection concerns the detection and mitigation of threats. 
% 
At the level of a collective, 
 these properties cannot generally be reduced to independent actions performed by its members. 
% 
Recovery, for example, 
 may require the redistribution of a failed component's responsibilities, 
 while optimization may require balancing competing local improvements 
 against their system-wide effects.

\emph{Self-organization} 
 denotes the process through which a system creates 
 and maintains an ordered structure without relying on an external controller. 
% 
In a \gls{CAS}, 
 this structure may consist of neighborhoods, hierarchies, spatial regions, coalitions, 
 or other patterns of interaction. 
% 
Self-organization is especially valuable when the population 
 and its environment change too frequently for a fixed organization to remain appropriate. 
% 
It is nevertheless a means rather than an objective: 
 the resulting structure is useful only insofar as it supports the intended collective behavior.

Self-organization is closely related to, but distinct from, \emph{emergence}.
%
Emergence describes the appearance of coherent macro-level properties 
 from interactions at the micro level, 
 where those properties cannot be directly assigned to any individual component. 
% 
The emergent outcome need not be beneficial or intended. 
%
The same feedback mechanisms that yield coordination may also produce 
oscillations, congestion, fragmentation, or other undesirable effects. 
%
Engineering a \gls{CAS} therefore requires both enabling desired emergent
 behavior and constraining the conditions under which harmful behavior can develop.

Adaptation also raises the question of stability. 
%
A system that reacts to every local variation 
 may continually reorganize  without converging to a useful state,
 while one that adapts too slowly may remain ineffective after 
 its operating conditions have changed. 
%
One relevant objective is \emph{self-stabilization}:
 after a transient perturbation, the collective should recover a state that satisfies 
 its intended properties without requiring global reinitialization.
%
Self-stabilizing coordination mechanisms provide a foundation for resilient \glspl{CAS}
 because they tolerate changes in local state, failures, and network topology
 while preserving system-level behavior~\cite{DBLP:journals/tomacs/ViroliABDP18}.

\section{Characteristic Properties and Engineering Implications}
\label{subsec:cas-properties}

Although \glspl{CAS} encompass different application domains, 
 several recurring properties shape how they can be engineered.
%
These include distribution, awareness, adaptation, interdependence, 
 collaboration, anonymity, scalability, and open-endedness~\cite{DBLP:journals/scp/AlrahmanNL20}.
%
Rather than occurring independently, these properties reinforce one another:
 local awareness supports adaptation, openness requires identity-independent
 interaction, and interdependence allows local actions to affect collective
 outcomes.

\subsection{Situatedness, Awareness, and Locality}
The entities of a \gls{CAS} operate within a logical or physical environment
 from which they obtain observations and on which their actions have effects.
%
Their knowledge can include information about their own run-time condition,
 such as available energy, computational load, or current role, as well as
 information about the surrounding context.
%
These two forms of knowledge are commonly referred to as
 \emph{self-awareness} and \emph{context-awareness}, respectively.
%
Awareness provides the information on which adaptation is based: an entity
 cannot respond meaningfully to a change that it cannot detect or estimate.
%
However, access to this information is normally local.
%
An entity observes only part of the relevant environmental state and
 communicates with a limited neighborhood.
%
Locality improves scalability and enables timely decisions, but prevents any
 component from assuming a continuously available global view.
%
Information about collective conditions must instead be estimated, propagated,
 or aggregated through the system.

\subsection{Scale, Openness, and Anonymity}
A \gls{CAS} may contain hundreds or thousands of entities, and its population
 is often open.
%
Components can join, leave, move, fail, or become temporarily unreachable,
 changing both the size of the system and its interaction topology.
%
Algorithms that depend on fixed membership, a stable ordering of entities,
 or global synchronization become increasingly fragile under these conditions.
%
For the same reason, interactions should not necessarily depend on the
 identities or permanent addresses of specific participants.
%
In many cases, an entity is relevant because of its current properties,
 capabilities, or role rather than because of its identity.
%
Anonymity in this context therefore denotes identity-independent interaction,
 not the concealment of identity for security purposes.
%
It enables groups to form dynamically around run-time conditions and allows
 components to be replaced without requiring the interaction logic to be
 reconfigured explicitly.

\subsection{Autonomy and Decentralized Control}
Individual entities retain control over at least part of their state and
 behavior.
%
Their decisions are based on local knowledge and may be made without
 authorization from a permanent central coordinator.
%
Decentralization removes single points of control, distributes computational
 load, and can improve responsiveness and fault tolerance.
%
At the same time, it makes global correctness harder to establish:
 locally reasonable decisions can conflict, while delayed or inconsistent
 information can prevent coordination or slow convergence.
%
Engineering mechanisms must therefore constrain local autonomy sufficiently
 to preserve collective properties without eliminating the ability of
 individual entities to adapt independently.

\subsection{Interdependence and Collaboration}
Entities in a \gls{CAS} are interdependent because their actions may alter
 conditions experienced by other entities.
%
This influence may be direct, as when components exchange messages, or
 indirect, as when an entity modifies a shared environment and thereby affects
 the observations or behavior of others.
%
Consequently, even an action selected from purely local information can have
 system-wide effects that were neither intended nor immediately observable.
%
At the same time, the intended behavior of the system often depends on
 collaboration: entities must combine their capabilities, observations, or
 actions to achieve objectives that cannot be attained individually.
%
Coordination mechanisms must therefore account for both deliberate cooperation
 and unintended interference mediated by the environment.

\subsection{Heterogeneity}
The members of a \gls{CAS} may differ in hardware, energy availability,
 communication capacity, data, mobility, and assigned role.
%
Heterogeneity can expand the capabilities of the collective, since different
 components contribute complementary resources and observations.
%
It also prevents uniform assumptions about what each component can compute,
 communicate, or learn.
%
Coordination mechanisms must therefore distinguish differences that are
 operationally relevant while avoiding unnecessary dependencies on the
 identities or ordering of individual entities.

\subsection{Resilience}
Failures and environmental changes are part of the expected operating
 conditions of a \gls{CAS} rather than exceptional events.
%
Resilience consequently does not mean preventing every disruption, but
 maintaining acceptable system-level behavior and recovering when the
 collective is perturbed.
%
Redundant components and decentralized control can support graceful
 degradation, although redundancy alone is insufficient:
 the collective must also detect changes and reorganize its activity around
 the resources that remain available.

\subsection{Engineering Implications}

These properties create an important engineering trade-off.
%
Local interaction, decentralization, openness, and identity-independent
 coordination enable scale and robustness, but also make the system difficult
 to predict, verify, and control.
%
Evaluation must consequently consider both levels of the system.
%
Component-level measures such as latency, energy consumption, or individual
 accuracy remain necessary, but they do not by themselves establish that the
 collective satisfies its intended function.
%
Macro-level properties such as coverage, convergence, fairness, stability,
 and resource utilization must be assessed as outcomes of the interactions
 among components.

\section{Relationship with Adjacent Paradigms}
\label{subsec:cas-related-paradigms}

\glspl{CAS} overlap with several established classes of systems,
 but are not interchangeable with them. 
%
\glspl{MAS} provide models and mechanisms
 for autonomy, interaction, coordination, and organization. 
% 
A \gls{CAS} can therefore often be modeled as a \gls{MAS}. 
%
The \gls{CAS} perspective, however, 
 places particular emphasis on how macroscopic behavior 
 and adaptation arise in potentially large, open, and dynamically structured populations.

Swarm systems represent a more specialized case. 
%
They typically comprise many relatively simple 
 and homogeneous agents that follow decentralized rules,
 producing robust and scalable behavior through redundancy 
 and local interaction~\cite{DBLP:journals/swarm/BrambillaFBD13}. 
% 
\glspl{CAS} retain these principles but also include heterogeneous entities, 
 explicit organizational roles, virtual services, and combinations of centralized 
 and decentralized mechanisms. 
% 
Accordingly, 
 an adaptive robotic swarm can be regarded as a specialized \gls{CAS}, 
 but a \gls{CAS} need not exhibit the uniformity 
 or bio-inspired organization commonly associated with swarms.

The distinction from self-adaptive systems concerns the locus of adaptation. 
%
A self-adaptive system may be a single software system 
 whose behavior is changed by a centralized adaptation manager. 
% 
In a \gls{CAS}, 
 adaptation is distributed across multiple interacting entities 
 and may also modify the relationships among them.
%
Likewise, 
 the \glsentryfull{IoT} and \glsentryfull{CPS} describe how computational entities are connected 
 to one another and to the physical world, but connectivity and physical embodiment do not alone imply 
 collective adaptation. 
%
They provide common substrates on which \glspl{CAS} can be realized.

\section{Representative Examples}
\label{sec:cas-examples}

The properties discussed above appear in systems with different components,
 interaction mechanisms, and objectives.
%
The following examples illustrate how the CAS perspective applies to physical cooperation, 
 shared-resource management, and distributed guidance.

\subsection{Heterogeneous Robotic Cooperation: Swarmanoid}
Swarmanoid is a robotic system comprising three types of robots 
 with complementary capabilities: flying \emph{eye-bots} explore the environment,
 ground-based \emph{foot-bots} provide mobility and transport,
 and \emph{hand-bots} climb structures and manipulate objects~\cite{DBLP:journals/ram/Dorigo13}.
%
In the demonstrated search-and-retrieval task,
 these capabilities are combined to locate and retrieve a book from a shelf.
%
Eye-bots establish an exploration and communication network,
 which foot-bots progressively complement on the ground,
 freeing flying robots for further exploration.
%
Foot-bots also assemble around hand-bots to transport them,
 then detach so that the hand-bots can climb and retrieve the object.
%
The organization of the system therefore changes as the task progresses:
 cooperation involves both information exchange 
 and the formation of temporary physical structures.
%
This example illustrates why collective decision making must account 
 for complementary capabilities and changing dependencies among agents.

\subsection{Shared Urban Resources: Bike-Sharing Systems}
Bike-sharing systems provide a socio-technical example 
 in which individual decisions continuously reshape the resources available
 to the population.
%
Loreti and Hillston~\cite{DBLP:conf/sfm/LoretiH16} use such a system as a running example
 for modeling and analyzing CASs with CARMA.
%
Users retrieve bicycles from stations and return them elsewhere,
 changing the spatial distribution of bicycles and empty docking spaces.
%
Consequently, 
 the feasibility of a trip depends on the accumulated effects of other users' choices.
%
The authors also identify prices and incentives as mechanisms
 through which users may adjust their behavior.
%
The collective objective is to maintain service availability,
 although individual users primarily pursue their own travel needs.
%
For cooperative learning and decision making,
 this distinction highlights the need to evaluate system-wide outcomes:
 improvements for individual participants do not necessarily produce
 a better allocation of shared resources.

\subsection{Distributed Evacuation Guidance: SAPERE}
The SAPERE framework includes an emergency evacuation example in a smart museum, 
 involving agents associated with visitors
 and security agents monitoring emergency and crowding conditions~\cite{DBLP:journals/taas/CastelliMRZ15}.
%
When an emergency is detected,
 agents associated with exits propagate information through the building,
 creating a field that visitor agents can follow toward an exit.
%
Aggregated information about crowd distribution influences the field,
 allowing guidance to account for congestion as well as proximity.
%
The evacuation service is thus sustained by interactions among monitoring, 
 information propagation, and visitor movement,
 rather than by an independently computed route for each visitor.
%
This example anticipates the field-based coordination mechanisms
 introduced in the next chapter:
 distributed observations are organized into a computational structure
 that supports collective action.
%
It also connects to the crowd-steering problems considered later in this thesis, 
 where local movement decisions must be assessed through their effects on the distribution of the population.

\section{Engineering Perspective Adopted in This Thesis}
\label{subsec:cas-thesis-perspective}

The systems considered in this thesis are cooperative many-agent \glspl{CAS}. 
%
Their members operate with partial observations, 
 communicate primarily with a time-varying local neighborhood, 
 and contribute to system-level objectives while retaining local state, data, and resource constraints. 
% 
Cooperation does not require every entity to perform the same action or maintain the same model.
%
Rather, 
 it requires mechanisms through which locally available information 
 and capabilities can be organized into behavior that benefits the collective.

Within this scope, 
 learning is one mechanism of adaptation, 
 while organization and deployment constitute two others. 
% 
Models and policies adapt from data and experience; 
 cooperation structures adapt to determine which entities exchange knowledge; 
 and deployment decisions adapt the placement of computation to the available edge--cloud resources. 
% 
These processes cannot be designed in isolation, 
 because each changes the context observed by the others. 
% 
The central engineering problem is therefore to express and control their interaction 
 at the level of the collective while preserving decentralized execution.

This perspective motivates the two complementary foundations developed in the following chapters. 
%
Macroprogramming provides abstractions for specifying collective behavior 
 without prescribing every local interaction individually~\cite{DBLP:journals/csur/Casadei23}, 
 whereas machine learning provides mechanisms for deriving models 
 and decisions when the relevant conditions cannot be exhaustively specified at design time. 
% 
Their integration supports \glspl{CAS} in which learning is not an isolated function executed 
 by independent devices, but a coordinated and adaptive part of the collective behavior.